\documentclass[ukrainian]{ucuproposal}


\projecttitle{Наскільки можна довіряти монокулярній глибині}
\projecttagline{Оцінка zero-shot моделей метричної глибини як єдиного джерела сприйняття для обльоту перешкод дроном DJI Tello}

\teammember{Мацелюх Максим}
\teammember{Михасяк Ярема}
\teammember{Пелешко Марко-Зенон}

\mentor{}

\track{research}

\repo{https://github.com/MakcDutch/AI-project} 

\proposaldate{2026-09-20}

\begin{document}
\maketitlepage

\section{Проблема та мотивація}
Дрон DJI~Tello не має ні лідара, ні стереопари, ні далекоміра: єдине
джерело інформації про геометрію сцени~--- одна RGB-камера. Тому будь-який
автономний обліт перешкод на цій платформі повністю спирається на те,
наскільки точно нейромережа вгадає відстані з одного кадру. Сучасні
foundation-моделі глибини заявляють \emph{метричну} оцінку zero-shot, тобто
у метрах і без донавчання під конкретну сцену~\cite{lin2025depthanything3,
yin2023metric3d}. Проте ці показники отримані на якісних знімках зі
стандартних камер, а бортовий кадр Tello~--- це стиснене відео 720p по
Wi-Fi, з вузьким полем зору, змазуванням під час поворотів і артефактами
кодування. Чи зберігається заявлена метрична точність у таких умовах і чи
її достатньо, щоб дрон не врізався,~--- питання відкрите, і саме воно є
предметом проєкту. Практична вага очевидна: якщо монокулярної глибини
достатньо, дешеві дрони без далекомірів стають придатними для
indoor-задач; якщо ні, важливо знати, де саме проходить межа. Наукова вага
в тому, що ми міряємо розрив між benchmark-умовами й реальною дешевою
камерою, а не повторюємо відомий результат.

\section{Огляд джерел та наявних рішень}
Лінія робіт Depth~Anything показала, що великий попередньо навчений
енкодер дає стійку відносну глибину майже в будь-якому домені, а окремі
metric-варіанти повертають відстані у метрах~\cite{yang2024depthanythingv2,
lin2025depthanything3}. Metric3D формулює завдання zero-shot метричної
оцінки явно й показує, що масштаб критично залежить від параметрів
камери~\cite{yin2023metric3d}. Depth~Any~Camera прямо досліджує перенесення
між типами камер і демонструє суттєве падіння \(\delta_1\) при зміні поля
зору~\cite{guo2025depthanycamera}; це найближча до нас робота, але вона
розглядає fisheye і 360\textdegree{} камери, а не дешеву дронову оптику в
замкненому контурі керування. З боку робототехніки Loquercio зі
співавторами показали, що навчена політика здатна вести квадрокоптер на
високій швидкості у захаращеному середовищі~\cite{loquercio2021highspeed},
проте там використано стереокамеру й потужний бортовий комп'ютер. Наша
ніша~--- перевірка найдешевшої можливої конфігурації: одна камера, жодних
обчислень на борту, готові моделі без донавчання.

\section{Дані}
Власний збір, три набори. (1)~\emph{Калібрувальний}: 30--50 кадрів
шахівниці для оцінки внутрішніх параметрів камери Tello через
\texttt{cv2.calibrateCamera}. Фокусна відстань у пікселях обов'язкова, бо
метричний вихід сучасних моделей масштабується саме через неї.
(2)~\emph{Валідаційний}: 200--400 кадрів у приміщеннях УКУ, для кожного
3--5 опорних точок із відстанями, виміряними лазерним далекоміром або
рулеткою. Навмисно включаємо складні поверхні: однотонні стіни, скло,
дзеркала, тонкі об'єкти, контрове світло. Це наш ground truth.
(3)~\emph{Польотний}: 20--30 записів прольотів через смугу перешкод із
синхронними логами відео, RC-команд, телеметрії та таймстемпів. Дані
збираємо самостійно, публікуємо під CC~BY~4.0. Людей у кадр не знімаємо,
тому персональних даних набір не містить; ліцензії й обсяг фіксуємо в
\texttt{data/README.md}.

\section{Підхід і методи}
Система складається з трьох шарів, які оцінюються незалежно.
\emph{Перцепція}: прогін 2--3 моделей метричної глибини
(Depth~Anything~3~metric, Metric3D, за наявності ресурсів
Depth~Any~Camera) на валідаційному наборі; окремо перевіряємо гіпотезу,
що номінального масштабування за фокусною відстанню недостатньо і
потрібне перекалібрування масштабу під конкретну камеру за кількома
опорними вимірами. \emph{Керування}: свідомо простий реактивний
контролер без карти й планувальника~--- карта глибини ділиться на
сектори, обирається найвільніший, спрацьовує пороговий гальмівний закон.
Простота тут методологічна: вона ізолює перцепцію як єдину змінну, що
досліджується. \emph{Аналіз відмов}: розбивка похибки за типами поверхонь
і зіставлення її з місцями реальних зіткнень. Обчислення виконуються на
ноутбуці, дрон передає відео і приймає команди по Wi-Fi.

\section{Оцінювання результатів}
Перцепція: AbsRel, RMSE та \(\delta_1\) проти виміряних відстаней, з
обов'язковою розбивкою за типами поверхонь; додатково затримка інференсу
та FPS. Baseline: (a)~оцінка глибини з оптичного потоку при відомому
русі дрона, (b)~модель відносної глибини з наївним масштабуванням за
номінальною фокусною. Керування: частота зіткнень на 10 прольотів,
success rate проходження смуги, мінімальна дистанція до перешкоди,
надлишковість траєкторії відносно прямої. Ключовий очікуваний
артефакт~--- залежність частоти зіткнень від точності глибини, тобто
емпіричний поріг придатності. Кількісно очікуємо \(\delta_1 \geq 0{,}8\)
на текстурованих поверхнях і суттєве падіння на однотонних і дзеркальних;
підтвердження або спростування цієї гіпотези є результатом у будь-якому
випадку.

\section{План і ризики}
\begin{keyfacts}
\textbf{Ключові цифри:} 3 моделі глибини \textperiodcentered\ 2 baseline
\textperiodcentered\ 200--400 кадрів з ground truth \textperiodcentered\ 5 тижнів до midterm
\end{keyfacts}
\textbf{До midterm (5 тижнів):} калібрування камери, пайплайн зчитування
відео, replay-стенд, збір валідаційного набору, перша таблиця похибок,
контролер у replay-режимі, перші живі прольоти. \textbf{До фіналу:}
повний цикл польотів, аналіз відмов, залежність зіткнень від точності,
абляції за частотою обробки й роздільною здатністю, постер.
\textbf{Ризики та план Б:} (1)~дрон недоступний або виходить з ладу~---
шар перцепції не потребує польотів узагалі й лишається самодостатнім
дослідженням; (2)~нестабільний Wi-Fi і затримка ламають живий контур~---
усі рішення оцінюються на replay-стенді, де записаний проліт
відтворюється детерміновано з керованою штучною затримкою;
(3)~бракує GPU для моделей у реальному часі~--- переходимо на малі
варіанти й обробку кожного \(N\)-го кадру, причому залежність якості від
частоти обробки сама стає частиною дослідження; (4)~замало часу на
польоти~--- скорочуємо смугу перешкод до фіксованого набору сценаріїв,
зберігаючи порівнюваність між моделями.

\section{Contribution statement}
\begin{contributions}
  \contribution{Мацелюх Максим}{калібрування камери, збір і розмітка валідаційного набору, вимірювання ground truth}
  \contribution{Михасяк Ярема}{прогін і порівняння моделей глибини, метрики точності, аналіз відмов за типами поверхонь}
  \contribution{Пелешко Марко-Зенон}{replay-стенд, реактивний контролер, польотні експерименти, репозиторій і постер}
\end{contributions}

\ucuprintbibliography

\end{document}
